\documentclass[10pt,a4paper]{amsart}
\usepackage{amsmath,amssymb,amsthm,graphicx}
\usepackage{verbatim}
\newenvironment{reference}{\comment}{\endcomment}
%\newenvironment{reference}{}{}
\newenvironment{slogan}{\comment}{\endcomment}
\newenvironment{history}{\comment}{\endcomment}

% For commutative diagrams we use Xy-pic
\usepackage[all]{xy}

% We use 2cell for 2-commutative diagrams.
\xyoption{2cell}
% Theorem environments.
%
\theoremstyle{plain}
\newtheorem{theorem}[subsection]{Theorem}
\newtheorem{proposition}[subsection]{Proposition}
\newtheorem{lemma}[subsection]{Lemma}

\theoremstyle{definition}
\newtheorem{definition}[subsection]{Definition}
\newtheorem{example}[subsection]{Example}
\newtheorem{exercise}[subsection]{Exercise}
\newtheorem{situation}[subsection]{Situation}

\theoremstyle{remark}
\newtheorem{remark}[subsection]{Remark}
\newtheorem{remarks}[subsection]{Remarks}

\numberwithin{equation}{subsection}

% Macros
%
\def\lim{\mathop{\mathrm{lim}}\nolimits}
\def\colim{\mathop{\mathrm{colim}}\nolimits}
\def\Spec{\mathop{\mathrm{Spec}}}
\def\Hom{\mathop{\mathrm{Hom}}\nolimits}
\def\Ext{\mathop{\mathrm{Ext}}\nolimits}
\def\SheafHom{\mathop{\mathcal{H}\!\mathit{om}}\nolimits}
\def\SheafExt{\mathop{\mathcal{E}\!\mathit{xt}}\nolimits}
\def\Sch{\mathit{Sch}}
\def\Mor{\mathop{\mathrm{Mor}}\nolimits}
\def\Ob{\mathop{\mathrm{Ob}}\nolimits}
\def\Sh{\mathop{\mathit{Sh}}\nolimits}
\def\NL{\mathop{N\!L}\nolimits}
\def\CH{\mathop{\mathrm{CH}}\nolimits}
\def\proetale{{pro\text{-}\acute{e}tale}}
\def\etale{{\acute{e}tale}}
\def\QCoh{\mathit{QCoh}}
\def\Ker{\mathop{\mathrm{Ker}}}
\def\Im{\mathop{\mathrm{Im}}}
\def\Coker{\mathop{\mathrm{Coker}}}
\def\Coim{\mathop{\mathrm{Coim}}}

% Boxtimes
%
\DeclareMathSymbol{\boxtimes}{\mathbin}{AMSa}{"02}

%
% Macros for moduli stacks/spaces
%
\def\QCohstack{\mathcal{QC}\!\mathit{oh}}
\def\Cohstack{\mathcal{C}\!\mathit{oh}}
\def\Spacesstack{\mathcal{S}\!\mathit{paces}}
\def\Quotfunctor{\mathrm{Quot}}
\def\Hilbfunctor{\mathrm{Hilb}}
\def\Curvesstack{\mathcal{C}\!\mathit{urves}}
\def\Polarizedstack{\mathcal{P}\!\mathit{olarized}}
\def\Complexesstack{\mathcal{C}\!\mathit{omplexes}}
% \Pic is the operator that assigns to X its picard group, usage \Pic(X)
% \Picardstack_{X/B} denotes the Picard stack of X over B
% \Picardfunctor_{X/B} denotes the Picard functor of X over B
\def\Pic{\mathop{\mathrm{Pic}}\nolimits}
\def\Picardstack{\mathcal{P}\!\mathit{ic}}
\def\Picardfunctor{\mathrm{Pic}}
\def\Deformationcategory{\mathcal{D}\!\mathit{ef}}
\usepackage{cleveref}
\crefname{originalreference}{original reference}{original references}
\Crefname{originalreference}{Original reference}{Original references}
\makeatletter
\newlabel{example-glue-points}{{0C1J}{source}{}{sourcepointer-1}{}}
\newlabel{example-glue-points@cref}{{[originalreference][0][]0C1J}{[1][0][]source}{}{}{}}
\makeatother
\crefname{originalreference}{original reference}{original references}
\Crefname{originalreference}{Original reference}{Original references}
\makeatletter
\newlabel{lemma-factor-almost-isomorphism}{{0C1L}{source}{}{sourcepointer-2}{}}
\newlabel{lemma-factor-almost-isomorphism@cref}{{[originalreference][0][]0C1L}{[1][0][]source}{}{}{}}
\makeatother
\crefname{originalreference}{original reference}{original references}
\Crefname{originalreference}{Original reference}{Original references}
\makeatletter
\newlabel{lemma-multicross-algebra}{{0C1V}{source}{}{sourcepointer-3}{}}
\newlabel{lemma-multicross-algebra@cref}{{[originalreference][0][]0C1V}{[1][0][]source}{}{}{}}
\makeatother
\title{Characterizations of multicross singularities by the delta-invariant and by iterated glueings}
\author{The Stacks Project Authors}
\date{stacks-project:curves.tex:local-snapshot-20261012; GNU FDL 1.2 or later}
\begin{document}
\maketitle
\noindent\textit{Editorial note.} All delivered mathematical text is the author's own, copied from the pinned Stacks Project chapter file curves.tex; no statement or proof step was rewritten, reordered or removed.\par\medskip
\noindent\textit{Editorial note.} The counted proof is the author's complete argument: the imported algebraic equivalence, the completed-ring inclusion with the branches/delta additivity computation, the normalization factorization, and the squishing exclusion by contradiction are all retained verbatim.\par\medskip
\noindent\textit{Editorial note.} The label lemma-multicross names the counted statement, and the proof environment immediately following it in the source is delivered verbatim as its proof.\par\medskip
\noindent\textit{Editorial note.} The counted statement is carried by the author's own lemma environment in the chapter source.\par\medskip
\noindent\textit{Allowed standard prerequisites (editorial).}\begin{itemize}
\item Multicross versus delta-invariant. Same chapter, Lemma lemma-multicross-algebra. Supplies the equivalence between the defining completed normal form and the delta-invariant equal to branches minus one.
\item Glueing of two points and the almost-isomorphism factorization. Same chapter, Example example-glue-points and Lemma lemma-factor-almost-isomorphism. Defines the glueing operation and factors the normalization map into glueings and tangent-vector squishings.
\end{itemize}\medskip
\ifdefined\hypertarget\hypertarget{sourcepointer-1}{}\fi
\noindent\textit{Source reference.} Original reference 0C1J: sources/stacks-curves/pinned-original-source.tex, source lines 3202--3202. Prerequisite pointer invoked at proof line(s) 3489; it names a standard fact of the same chapter that the counted argument uses, and every step of the counted proof built on it is retained verbatim in the delivered body.\par\medskip
\ifdefined\hypertarget\hypertarget{sourcepointer-2}{}\fi
\noindent\textit{Source reference.} Original reference 0C1L: sources/stacks-curves/pinned-original-source.tex, source lines 3257--3257. Prerequisite pointer invoked at proof line(s) 3515; it names a standard fact of the same chapter that the counted argument uses, and every step of the counted proof built on it is retained verbatim in the delivered body.\par\medskip
\ifdefined\hypertarget\hypertarget{sourcepointer-3}{}\fi
\noindent\textit{Source reference.} Original reference 0C1V: sources/stacks-curves/pinned-original-source.tex, source lines 3431--3431. Prerequisite pointer invoked at proof line(s) 3494; it names a standard fact of the same chapter that the counted argument uses, and every step of the counted proof built on it is retained verbatim in the delivered body.\par\medskip
\begin{lemma}
\label{lemma-multicross}
Let $k$ be an algebraically closed field. Let $X$ be a reduced algebraic
$1$-dimensional $k$-scheme. Let $x \in X$. The following are equivalent
\begin{enumerate}
\item $x$ defines a multicross singularity,
\item the $\delta$-invariant of $X$ at $x$ is the
number of branches of $X$ at $x$ minus $1$,
\item there is a sequence of morphisms
$U_n \to U_{n - 1} \to \ldots \to U_0 = U \subset X$
where $U$ is an open neighbourhood of $x$, where
$U_n$ is nonsingular, and where each $U_i \to U_{i - 1}$
is the glueing of two points as in Example \ref{example-glue-points}.
\end{enumerate}
\end{lemma}

\begin{proof}
The equivalence of (1) and (2) is Lemma \ref{lemma-multicross-algebra}.

\medskip\noindent
Assume (3). We will argue by descending induction on $i$ that all singularities
of $U_i$ are multicross. This is true for $U_n$ as $U_n$ has no singular points.
If $U_i$ is gotten from $U_{i + 1}$ by glueing $a, b \in U_{i + 1}$
to a point $c \in U_i$, then we see that
$$
\mathcal{O}_{U_i, c}^\wedge \subset
\mathcal{O}_{U_{i + 1}, a}^\wedge \times \mathcal{O}_{U_{i + 1}, b}^\wedge
$$
is the set of elements having the same residue classes in $k$.
Thus the number of branches at $c$ is the sum of the number of
branches at $a$ and $b$, and the $\delta$-invariant at $c$
is the sum of the $\delta$-invariants at $a$ and $b$ plus $1$
(because the displayed inclusion has codimension $1$).
This proves that (2) holds as desired.

\medskip\noindent
Assume the equivalent conditions (1) and (2). We may choose an open
$U \subset X$ such that $x$ is the only singular point of $U$.
Then we apply Lemma \ref{lemma-factor-almost-isomorphism} to
the normalization morphism
$$
U^\nu = U_n \to U_{n - 1} \to \ldots \to U_1 \to U_0 = U
$$
All we have to do is show that in none of the steps we are
squishing a tangent vector. Suppose $U_{i + 1} \to U_i$ is the
smallest $i$ such that this is the squishing of a tangent
vector $\theta$ at $u' \in U_{i + 1}$ lying over $u \in U_i$.
Arguing as above, we see that $u_i$ is a multicross singularity
(because the maps $U_i \to \ldots \to U_0$ are glueing of
pairs of points). But now the number of branches at $u'$ and $u$
is the same and the $\delta$-invariant of $U_i$ at $u$
is $1$ bigger than the $\delta$-invariant of $U_{i + 1}$ at $u'$.
By Lemma \ref{lemma-multicross-algebra}
this implies that $u$ cannot be a multicross singularity which
is a contradiction.
\end{proof}

\end{document}
